\subsection{Erzeugung durch Adjunktion}

Ein Alphabet \(A\) ist eine beliebige Menge. Wir beginnen mit der leeren
Menge und fügen einen Buchstaben zusammen mit der Anzahl der bisher
gebildeten Einträge an. So entstehen beispielsweise
\[
  \varnothing,\qquad \{(0,a)\},\qquad
  \{(0,a),(1,b)\},\qquad \{(0,a),(1,b),(2,a)\}.
\]
Die Paare halten Reihenfolge und Wiederholungen fest. Der Aufbau verwendet
nur endliche Teilmengen, ihre Kardinalzahl und die Adjunktion aus Band~20.
Wie bei der Konstruktion der natürlichen Zahlen in Band~10 nehmen wir
anschließend die kleinste unter den angegebenen Regeln abgeschlossene
Menge.

\newcommand{\FiniteOccurrenceSets}[1]{\mathcal E_{#1}}
\newcommand{\OccurrenceAdjunction}[2]{J(#1,#2)}
\newcommand{\WordGenerationFamilies}[1]{\mathfrak C_{#1}}

\FormulaDefDeltaK[Umgebungsbereich der endlichen Teilmengen]{%
  \FiniteOccurrenceSets A\coloneqq\FinSubsets{\mathbb N\times A}
}{FiniteOccurrenceSetUniverseDef}{%
  \DeltaRow{Mengen}{A}
  \DeltaRow{\textbf{Neues Symbol}}{}
  \DeltaPrem{Umgebungsbereich}{\FiniteOccurrenceSets A}
}

Endlich sind die Elemente von \(\FiniteOccurrenceSets A\); der
Umgebungsbereich selbst braucht nicht endlich zu sein.

\Needspace{16\baselineskip}
\FormulaThmDeltaK[Elemente des Umgebungsbereichs]{%
  w\in\FiniteOccurrenceSets A\eqvdash
    w\subseteq\mathbb N\times A\land\Finite w
}{FiniteOccurrenceSetUniverseMembership}{\DeltaRow{Mengen}{A\dsep w}}
\begin{tabproofsplitwide}
  \Needspace{9\baselineskip}
  \proofpartwide{Vom Umgebungsbereich zur endlichen Teilmenge}
    \proofstepwidestar[1]{w\in\FiniteOccurrenceSets A}{\rA}
    \proofstepwidestar[]{\FiniteOccurrenceSets A=\FinSubsets{\mathbb N\times A}}{%
      \FormulaRefAuto{FiniteOccurrenceSetUniverseDef}[def]}
    \proofstepwidestar[1]{w\in\FinSubsets{\mathbb N\times A}}{\rIE{2,1}}
    \proofstepwidestar[1]{w\subseteq\mathbb N\times A\land\Finite w}{%
      \FormulaRefAuto{F\in\FinSubsets{M}\eqvdash F\subseteq M\land\Finite{F}}[thm]{3}}
  \closeproofpartwide
  \Needspace{9\baselineskip}
  \proofpartwide{Von der endlichen Teilmenge zum Umgebungsbereich}
    \proofstepwidestar[1]{w\subseteq\mathbb N\times A\land\Finite w}{\rA}
    \proofstepwidestar[1]{w\in\FinSubsets{\mathbb N\times A}}{%
      \FormulaRefAuto{F\in\FinSubsets{M}\eqvdash F\subseteq M\land\Finite{F}}[thm]{1}}
    \proofstepwidestar[]{\FiniteOccurrenceSets A=\FinSubsets{\mathbb N\times A}}{%
      \FormulaRefAuto{FiniteOccurrenceSetUniverseDef}[def]}
    \proofstepwidestar[1]{w\in\FiniteOccurrenceSets A}{\rIE{3,2}}
  \closeproofpartwide
\end{tabproofsplitwide}

\FormulaDefDeltaK[Adjunktion an der Kardinalzahlposition]{%
  w\in\FiniteOccurrenceSets A\dsep a\in A\vdash
    \OccurrenceAdjunction w a\coloneqq w\cup\{(\FiniteCard(w),a)\}
}{OccurrenceCardinalAdjunctionDef}{%
  \DeltaRow{Mengen}{A\dsep w\dsep a}
  \DeltaRow{\textbf{Neues Symbol}}{}
  \DeltaPrem{Erzeugungsschritt}{\OccurrenceAdjunction w a}
}

Die Kardinalzahl \(\FiniteCard(w)\) ist bereits aus Band~20 bekannt. Sie zählt
hier die Elemente einer endlichen Menge; eine Wortlänge wird für den
Erzeugungsschritt noch nicht vorausgesetzt. Auf einer beliebigen endlichen
Teilmenge muss das angefügte Paar nicht neu sein. Für die Abgeschlossenheit
des Umgebungsbereichs genügt deshalb die Adjunktion eines beliebigen
Elements.

\Needspace{16\baselineskip}
\FormulaThmDeltaK[Die leere Menge liegt im Umgebungsbereich]{%
  \varnothing\in\FiniteOccurrenceSets A
}{FiniteOccurrenceUniverseEmpty}{\DeltaRow{Mengen}{A}}
\begin{tabproofwide}
  \proofstepwidestar[]{\varnothing\subseteq\mathbb N\times A}{%
    \FormulaRefAuto{\varnothing\subseteq A}[thm]}
  \proofstepwidestar[]{\Finite\varnothing}{\FormulaRefAuto{FiniteEmptyAxiom}[ax]}
  \proofstepwidestar[]{\varnothing\subseteq\mathbb N\times A\land\Finite\varnothing}{\rAI{1,2}}
  \proofstepwidestar[]{\varnothing\in\FiniteOccurrenceSets A}{%
    \FormulaRefAuto{FiniteOccurrenceSetUniverseMembership}[thm]{3}}
\end{tabproofwide}

\Needspace{16\baselineskip}
\FormulaThmDeltaK[Der Erzeugungsschritt bleibt im Umgebungsbereich]{%
  w\in\FiniteOccurrenceSets A\dsep a\in A\vdash
    \OccurrenceAdjunction w a\in\FiniteOccurrenceSets A
}{FiniteOccurrenceUniverseClosed}{\DeltaRow{Mengen}{A\dsep w\dsep a}}
\begin{tabproofwide}
  \proofstepwidestar[1]{w\in\FiniteOccurrenceSets A}{\rA}
  \proofstepwidestar[2]{a\in A}{\rA}
  \proofstepwidestar[1]{w\subseteq\mathbb N\times A\land\Finite w}{%
    \FormulaRefAuto{FiniteOccurrenceSetUniverseMembership}[thm]{1}}
  \proofstepwidestar[1]{w\subseteq\mathbb N\times A}{\rAEa{3}}
  \proofstepwidestar[1]{\Finite w}{\rAEb{3}}
  \proofstepwidestar[1]{\FiniteCard(w)\in\mathbb N}{\FormulaRefAuto{FiniteCardNatural}[thm]{5}}
  \proofstepwidestar[1,2]{(\FiniteCard(w),a)\in\mathbb N\times A}{%
    \FormulaRefAuto{(a,b)\in A\times B\eqvdash a\in A\land b\in B}[thm]{\rAI{6,2}}}
  \proofstepwidestar[1,2]{\{(\FiniteCard(w),a)\}\subseteq\mathbb N\times A}{%
    \FormulaRefAuto{a\in A\vdash\{a\}\subseteq A}[thm]{7}}
  \proofstepwidestar[1,2]{w\cup\{(\FiniteCard(w),a)\}\subseteq\mathbb N\times A}{%
    \FormulaRefAuto{A\subseteq C\dsep B\subseteq C\vdash A\cup B\subseteq C}[thm]{4,8}}
  \proofstepwidestar[1]{\Finite{w\cup\{(\FiniteCard(w),a)\}}}{%
    \FormulaRefAuto{FiniteInsert}[thm]{5}}
  \proofstepwidestar[1,2]{w\cup\{(\FiniteCard(w),a)\}\in\FiniteOccurrenceSets A}{%
    \FormulaRefAuto{FiniteOccurrenceSetUniverseMembership}[thm]{\rAI{9,10}}}
  \proofstepwidestar[1,2]{\OccurrenceAdjunction w a=w\cup\{(\FiniteCard(w),a)\}}{%
    \FormulaRefAuto{OccurrenceCardinalAdjunctionDef}[def]{1,2}}
  \proofstepwidestar[1,2]{\OccurrenceAdjunction w a\in\FiniteOccurrenceSets A}{\rIE{12,11}}
\end{tabproofwide}

\subsection{Die kleinste abgeschlossene Menge}

\noindent\begin{minipage}{\linewidth}
\FormulaDefDeltaK[Familie der abgeschlossenen Mengen]{%
  \begin{aligned}[t]
    \WordGenerationFamilies A\coloneqq\bigl\{C\in\powerset(\FiniteOccurrenceSets A)\bigm|{}&
      \varnothing\in C\ \land{}\\[-2pt]
    &\forall w\in C\,\forall a\in A\,
      (\OccurrenceAdjunction w a\in C)\bigr\}
  \end{aligned}
}{WordGenerationFamilyDef}{%
  \DeltaRow{Mengen}{A\dsep C\dsep w\dsep a}
  \DeltaRow{\textbf{Neues Symbol}}{}
  \DeltaPrem{Abgeschlossene Mengen}{\WordGenerationFamilies A}
}
\end{minipage}\par

\Needspace{18\baselineskip}
\noindent\begin{minipage}{\linewidth}
\FormulaThmDeltaK[Elementkriterium der abgeschlossenen Mengen]{%
  \begin{aligned}[t]
    C\in\WordGenerationFamilies A\eqvdash{}&
    C\subseteq\FiniteOccurrenceSets A\ \land{}\\[-2pt]
    &\varnothing\in C\ \land
      \forall w\in C\,\forall a\in A\,(\OccurrenceAdjunction w a\in C)
  \end{aligned}
}{WordGenerationFamilyMembership}{\DeltaRow{Mengen}{A\dsep C\dsep w\dsep a}}
\end{minipage}\par
\begin{tabproofsplitwide}
  \Needspace{9\baselineskip}
  \proofpartwide{Eigenschaften eines Familienglieds}
    \proofstepwidestar[1]{C\in\WordGenerationFamilies A}{\rA}
    \proofstepwidestar[1]{\begin{aligned}[t]
      &C\in\powerset(\FiniteOccurrenceSets A)\land{}\\[-2pt]
      &\varnothing\in C\land
        \forall w\in C\,\forall a\in A\,(\OccurrenceAdjunction w a\in C)
    \end{aligned}}{%
      \FormulaRefAuto{x\in\{u\in A\mid P(u)\}\eqvdash x\in A\land P(x)}[thm]{%
        \rIE{\FormulaRefAuto{WordGenerationFamilyDef}[def],1}}}
    \proofstepwidestar[1]{C\subseteq\FiniteOccurrenceSets A}{%
      \FormulaRefAuto{x\in\powerset(A)\eqvdash x\subseteq A}[thm]{\rAEa{2}}}
    \proofstepwidestar[1]{\begin{aligned}[t]
      &C\subseteq\FiniteOccurrenceSets A\land{}\\[-2pt]
      &\varnothing\in C\land
        \forall w\in C\,\forall a\in A\,(\OccurrenceAdjunction w a\in C)
    \end{aligned}}{\rAI{3,\rAEb{2}}}
  \closeproofpartwide
  \Needspace{9\baselineskip}
  \proofpartwide{Aufnahme einer abgeschlossenen Menge}
    \proofstepwidestar[1]{\begin{aligned}[t]
      &C\subseteq\FiniteOccurrenceSets A\land{}\\[-2pt]
      &\varnothing\in C\land
        \forall w\in C\,\forall a\in A\,(\OccurrenceAdjunction w a\in C)
    \end{aligned}}{\rA}
    \proofstepwidestar[1]{C\in\powerset(\FiniteOccurrenceSets A)}{%
      \FormulaRefAuto{x\in\powerset(A)\eqvdash x\subseteq A}[thm]{\rAEa{1}}}
    \proofstepwidestar[1]{\begin{aligned}[t]
      &C\in\powerset(\FiniteOccurrenceSets A)\land{}\\[-2pt]
      &\varnothing\in C\land
        \forall w\in C\,\forall a\in A\,(\OccurrenceAdjunction w a\in C)
    \end{aligned}}{\rAI{2,\rAEb{1}}}
    \proofstepwidestar[1]{C\in\WordGenerationFamilies A}{%
      \rIE{\FormulaRefAuto{WordGenerationFamilyDef}[def],%
        \FormulaRefAuto{x\in\{u\in A\mid P(u)\}\eqvdash x\in A\land P(x)}[thm]{3}}}
  \closeproofpartwide
\end{tabproofsplitwide}

\Needspace{16\baselineskip}
\FormulaThmDeltaK[Der Umgebungsbereich ist eine abgeschlossene Menge]{%
  \FiniteOccurrenceSets A\in\WordGenerationFamilies A
}{WordGenerationUniverseMember}{\DeltaRow{Mengen}{A\dsep w\dsep a}}
\begin{tabproofwide}
  \proofstepwidestar[]{\FiniteOccurrenceSets A\subseteq\FiniteOccurrenceSets A}{%
    \FormulaRefAuto{A\subseteq A}[thm]}
  \proofstepwidestar[]{\varnothing\in\FiniteOccurrenceSets A}{%
    \FormulaRefAuto{FiniteOccurrenceUniverseEmpty}[thm]}
  \proofstepwidestar[3]{w\in\FiniteOccurrenceSets A}{\rA}
  \proofstepwidestar[4]{a\in A}{\rA}
  \proofstepwidestar[3,4]{\OccurrenceAdjunction w a\in\FiniteOccurrenceSets A}{%
    \FormulaRefAuto{FiniteOccurrenceUniverseClosed}[thm]{3,4}}
  \proofstepwidestar[]{\forall w\in\FiniteOccurrenceSets A\,\forall a\in A\,
    (\OccurrenceAdjunction w a\in\FiniteOccurrenceSets A)}{\rUI{\rRI{3,\rUI{\rRI{4,5}}}}}
  \proofstepwidestar[]{\FiniteOccurrenceSets A\in\WordGenerationFamilies A}{%
    \FormulaRefAuto{WordGenerationFamilyMembership}[thm]{\rAI{1,\rAI{2,6}}}}
\end{tabproofwide}

\Needspace{16\baselineskip}
\FormulaThmDeltaK[Die Familie abgeschlossener Mengen ist nicht leer]{%
  \WordGenerationFamilies A\neq\varnothing
}{WordGenerationFamilyNonempty}{\DeltaRow{Mengen}{A}}
\begin{tabproofwide}
  \proofstepwidestar[]{\FiniteOccurrenceSets A\in\WordGenerationFamilies A}{%
    \FormulaRefAuto{WordGenerationUniverseMember}[thm]}
  \proofstepwidestar[]{\exists C\,(C\in\WordGenerationFamilies A)}{\rEI{1}}
  \proofstepwidestar[]{\WordGenerationFamilies A\neq\varnothing}{%
    \FormulaRefAuto{A\neq\varnothing\eqvdash\exists x\,(x\in A)}[thm]{2}}
\end{tabproofwide}

Damit ist der folgende Durchschnitt nach Band~03 definiert.

\FormulaDefDeltaK[Menge der endlichen Wörter]{%
  \FiniteWords A\coloneqq\bigcap\WordGenerationFamilies A
}{FiniteWordSetDef}{%
  \DeltaRow{Mengen}{A}
  \DeltaRow{\textbf{Neues Symbol}}{}
  \DeltaPrem{Endliche Wörter}{\FiniteWords A}
}

Die Elemente von \(\FiniteWords A\) heißen \textbf{endliche Wörter über
\(A\)}. Die Definition fordert ausschließlich den leeren Ausgangsgegenstand
und den Erzeugungsschritt. Die folgenden Sätze zeigen, dass ihr
Durchschnitt genau die kleinste abgeschlossene Menge ist.

\Needspace{16\baselineskip}
\FormulaThmDeltaK[Elementkriterium des erzeugten Wortbereichs]{%
  w\in\FiniteWords A\leftrightarrow
    \forall C\in\WordGenerationFamilies A\,(w\in C)
}{GeneratedWordMembership}{\DeltaRow{Mengen}{A\dsep w\dsep C}}
\begin{tabproofwide}
  \proofstepwidestar[]{\WordGenerationFamilies A\neq\varnothing}{%
    \FormulaRefAuto{WordGenerationFamilyNonempty}[thm]}
  \proofstepwidestar[]{\begin{aligned}[t]
    &w\in\bigcap\WordGenerationFamilies A\leftrightarrow{}\\[-2pt]
    &\forall C\in\WordGenerationFamilies A\,(w\in C)
  \end{aligned}}{%
    \FormulaRefAuto{M\neq\varnothing\vdash x\in\bigcap M\leftrightarrow\forall X\in M\,(x\in X)}[thm]{1}}
  \proofstepwidestar[]{\FiniteWords A=\bigcap\WordGenerationFamilies A}{%
    \FormulaRefAuto{FiniteWordSetDef}[def]}
  \proofstepwidestar[]{\begin{aligned}[t]
    &w\in\FiniteWords A\leftrightarrow{}\\[-2pt]
    &\forall C\in\WordGenerationFamilies A\,(w\in C)
  \end{aligned}}{\rIE{3,2}}
\end{tabproofwide}

\Needspace{16\baselineskip}
\FormulaThmDeltaK[Jede abgeschlossene Menge enthält alle Wörter]{%
  C\in\WordGenerationFamilies A\vdash\FiniteWords A\subseteq C
}{FiniteWordSetGenerationLeast}{\DeltaRow{Mengen}{A\dsep C\dsep w\dsep D}}
\begin{tabproofwide}
  \proofstepwidestar[1]{C\in\WordGenerationFamilies A}{\rA}
  \proofstepwidestar[2]{w\in\FiniteWords A}{\rA}
  \proofstepwidestar[]{w\in\FiniteWords A\leftrightarrow
    \forall D\in\WordGenerationFamilies A\,(w\in D)}{%
    \FormulaRefAuto{GeneratedWordMembership}[thm]}
  \proofstepwidestar[2]{\forall D\in\WordGenerationFamilies A\,(w\in D)}{\rRE{2,\rLREa{3}}}
  \proofstepwidestar[1,2]{w\in C}{\rRE{1,\rUE{4}}}
  \proofstepwidestar[1]{\FiniteWords A\subseteq C}{%
    \FormulaRefAuto{SubsetIntroductionRule}[thm]{[2]\vdots 5}}
\end{tabproofwide}

\Needspace{16\baselineskip}
\FormulaThmDeltaK[Die Wortmenge liegt im Umgebungsbereich]{%
  \FiniteWords A\subseteq\FiniteOccurrenceSets A
}{FiniteWordSetSubsetUniverse}{\DeltaRow{Mengen}{A}}
\begin{tabproofwide}
  \proofstepwidestar[]{\FiniteOccurrenceSets A\in\WordGenerationFamilies A}{%
    \FormulaRefAuto{WordGenerationUniverseMember}[thm]}
  \proofstepwidestar[]{\FiniteWords A\subseteq\FiniteOccurrenceSets A}{%
    \FormulaRefAuto{FiniteWordSetGenerationLeast}[thm]{1}}
\end{tabproofwide}

\Needspace{16\baselineskip}
\FormulaThmDeltaK[Jedes Wort liegt im Umgebungsbereich]{%
  w\in\FiniteWords A\vdash w\in\FiniteOccurrenceSets A
}{FiniteWordInUniverse}{\DeltaRow{Mengen}{A\dsep w}}
\begin{tabproofwide}
  \proofstepwidestar[1]{w\in\FiniteWords A}{\rA}
  \proofstepwidestar[]{\FiniteWords A\subseteq\FiniteOccurrenceSets A}{%
    \FormulaRefAuto{FiniteWordSetSubsetUniverse}[thm]}
  \proofstepwidestar[1]{w\in\FiniteOccurrenceSets A}{%
    \FormulaRefAuto{A\subseteq B\dsep x\in A\vdash x\in B}[thm]{2,1}}
\end{tabproofwide}

\Needspace{16\baselineskip}
\FormulaThmDeltaK[Jedes Wort ist eine endliche Menge]{%
  w\in\FiniteWords A\vdash\Finite w
}{FiniteWordUnderlyingFinite}{\DeltaRow{Mengen}{A\dsep w}}
\begin{tabproofwide}
  \proofstepwidestar[1]{w\in\FiniteWords A}{\rA}
  \proofstepwidestar[1]{w\in\FiniteOccurrenceSets A}{\FormulaRefAuto{FiniteWordInUniverse}[thm]{1}}
  \proofstepwidestar[1]{w\subseteq\mathbb N\times A\land\Finite w}{%
    \FormulaRefAuto{FiniteOccurrenceSetUniverseMembership}[thm]{2}}
  \proofstepwidestar[1]{\Finite w}{\rAEb{3}}
\end{tabproofwide}

\Needspace{18\baselineskip}
\FormulaThmDeltaK[Das leere Wort]{%
  \varnothing\in\FiniteWords A
}{GeneratedWordEmpty}{\DeltaRow{Mengen}{A\dsep C\dsep w\dsep a}}
\begin{tabproofwide}
  \proofstepwidestar[1]{C\in\WordGenerationFamilies A}{\rA}
  \proofstepwidestar[1]{\begin{aligned}[t]
    &C\subseteq\FiniteOccurrenceSets A\land\varnothing\in C\land{}\\[-2pt]
    &\forall w\in C\,\forall a\in A\,(\OccurrenceAdjunction w a\in C)
  \end{aligned}}{\FormulaRefAuto{WordGenerationFamilyMembership}[thm]{1}}
  \proofstepwidestar[1]{\varnothing\in C}{\rAEa{\rAEb{2}}}
  \proofstepwidestar[]{\forall C\in\WordGenerationFamilies A\,(\varnothing\in C)}{\rUI{\rRI{1,3}}}
  \proofstepwidestar[]{\varnothing\in\FiniteWords A\leftrightarrow
    \forall C\in\WordGenerationFamilies A\,(\varnothing\in C)}{%
    \FormulaRefAuto{GeneratedWordMembership}[thm]}
  \proofstepwidestar[]{\varnothing\in\FiniteWords A}{\rRE{4,\rLREb{5}}}
\end{tabproofwide}

\Needspace{16\baselineskip}
\FormulaThmDeltaK[Wörter bleiben unter dem Erzeugungsschritt abgeschlossen]{%
  w\in\FiniteWords A\dsep a\in A\vdash\OccurrenceAdjunction w a\in\FiniteWords A
}{GeneratedWordClosed}{\DeltaRow{Mengen}{A\dsep w\dsep a\dsep C\dsep u\dsep b}}
\begin{tabproofwide}
  \proofstepwidestar[1]{w\in\FiniteWords A}{\rA}
  \proofstepwidestar[2]{a\in A}{\rA}
  \proofstepwidestar[1]{w\in\FiniteOccurrenceSets A}{\FormulaRefAuto{FiniteWordInUniverse}[thm]{1}}
  \proofstepwidestar[4]{C\in\WordGenerationFamilies A}{\rA}
  \proofstepwidestar[4]{\FiniteWords A\subseteq C}{\FormulaRefAuto{FiniteWordSetGenerationLeast}[thm]{4}}
  \proofstepwidestar[1,4]{w\in C}{%
    \FormulaRefAuto{A\subseteq B\dsep x\in A\vdash x\in B}[thm]{5,1}}
  \proofstepwidestar[4]{\begin{aligned}[t]
    &C\subseteq\FiniteOccurrenceSets A\land\varnothing\in C\land{}\\[-2pt]
    &\forall u\in C\,\forall b\in A\,(\OccurrenceAdjunction u b\in C)
  \end{aligned}}{\FormulaRefAuto{WordGenerationFamilyMembership}[thm]{4}}
  \proofstepwidestar[4]{\forall u\in C\,\forall b\in A\,(\OccurrenceAdjunction u b\in C)}{\rAEb{\rAEb{7}}}
  \proofstepwidestar[1,4]{\forall b\in A\,(\OccurrenceAdjunction w b\in C)}{\rRE{6,\rUE{8}}}
  \proofstepwidestar[1,2,4]{\OccurrenceAdjunction w a\in C}{\rRE{2,\rUE{9}}}
  \proofstepwidestar[1,2]{\forall C\in\WordGenerationFamilies A\,(\OccurrenceAdjunction w a\in C)}{\rUI{\rRI{4,10}}}
  \proofstepwidestar[1,2]{\begin{aligned}[t]
    &\OccurrenceAdjunction w a\in\FiniteWords A\leftrightarrow{}\\[-2pt]
    &\forall C\in\WordGenerationFamilies A\,(\OccurrenceAdjunction w a\in C)
  \end{aligned}}{\FormulaRefAuto{GeneratedWordMembership}[thm]}
  \proofstepwidestar[1,2]{\OccurrenceAdjunction w a\in\FiniteWords A}{\rRE{11,\rLREb{12}}}
\end{tabproofwide}

\Needspace{16\baselineskip}
\FormulaThmDeltaK[Die Wortmenge ist eine abgeschlossene Menge]{%
  \FiniteWords A\in\WordGenerationFamilies A
}{FiniteWordSetGenerationClosed}{\DeltaRow{Mengen}{A\dsep w\dsep a}}
\begin{tabproofwide}
  \proofstepwidestar[]{\FiniteWords A\subseteq\FiniteOccurrenceSets A}{\FormulaRefAuto{FiniteWordSetSubsetUniverse}[thm]}
  \proofstepwidestar[]{\varnothing\in\FiniteWords A}{\FormulaRefAuto{GeneratedWordEmpty}[thm]}
  \proofstepwidestar[3]{w\in\FiniteWords A}{\rA}
  \proofstepwidestar[4]{a\in A}{\rA}
  \proofstepwidestar[3,4]{\OccurrenceAdjunction w a\in\FiniteWords A}{\FormulaRefAuto{GeneratedWordClosed}[thm]{3,4}}
  \proofstepwidestar[]{\forall w\in\FiniteWords A\,\forall a\in A\,(\OccurrenceAdjunction w a\in\FiniteWords A)}{\rUI{\rRI{3,\rUI{\rRI{4,5}}}}}
  \proofstepwidestar[]{\FiniteWords A\in\WordGenerationFamilies A}{%
    \FormulaRefAuto{WordGenerationFamilyMembership}[thm]{\rAI{1,\rAI{2,6}}}}
\end{tabproofwide}

\Needspace{16\baselineskip}
\noindent\begin{minipage}{\linewidth}
\FormulaThmDeltaK[Minimalität der erzeugten Wortmenge]{%
  \begin{aligned}[t]
    &C\subseteq\FiniteOccurrenceSets A\dsep\varnothing\in C\dsep{}\\[-2pt]
    &\forall w\in C\,\forall a\in A\,(\OccurrenceAdjunction w a\in C)
      \vdash\FiniteWords A\subseteq C
  \end{aligned}
}{GeneratedWordMinimality}{\DeltaRow{Mengen}{A\dsep C\dsep w\dsep a}}
\end{minipage}\par
\begin{tabproofwide}
  \proofstepwidestar[1]{C\subseteq\FiniteOccurrenceSets A}{\rA}
  \proofstepwidestar[2]{\varnothing\in C}{\rA}
  \proofstepwidestar[3]{\forall w\in C\,\forall a\in A\,(\OccurrenceAdjunction w a\in C)}{\rA}
  \proofstepwidestar[1,2,3]{C\in\WordGenerationFamilies A}{%
    \FormulaRefAuto{WordGenerationFamilyMembership}[thm]{\rAI{1,\rAI{2,3}}}}
  \proofstepwidestar[1,2,3]{\FiniteWords A\subseteq C}{\FormulaRefAuto{FiniteWordSetGenerationLeast}[thm]{4}}
\end{tabproofwide}

\subsection{Induktion aus der Konstruktion}

Die Minimalität liefert bereits ein Induktionsprinzip für das konkrete
Modell. Es wird hier aus der Schnittkonstruktion bewiesen. Die späteren
Strukturaxiome werden dafür nicht verwendet.

\Needspace{16\baselineskip}
\noindent\begin{minipage}{\linewidth}
\FormulaThmDeltaK[Induktion für erzeugte Wörter]{%
  \begin{aligned}[t]
    &P(\varnothing)\dsep{}\\[-2pt]
    &\forall w\in\FiniteWords A\,\forall a\in A\,
      \bigl(P(w)\rightarrow P(\OccurrenceAdjunction w a)\bigr)\dsep{}\\[-2pt]
    &v\in\FiniteWords A\vdash P(v)
  \end{aligned}
}{GeneratedWordInduction}{%
  \DeltaRow{Mengen}{A\dsep v\dsep w\dsep a}
  \DeltaRow{Prädikat}{P}
  \DeltaRow{Abkürzung}{U=\{w\in\FiniteWords A\mid P(w)\}}
}
\end{minipage}\par
\begin{tabproofwide}
  \proofstepwidestar[1]{P(\varnothing)}{\rA}
  \proofstepwidestar[2]{\forall w\in\FiniteWords A\,\forall a\in A\,
    \bigl(P(w)\rightarrow P(\OccurrenceAdjunction w a)\bigr)}{\rA}
  \proofstepwidestar[3]{v\in\FiniteWords A}{\rA}
  \proofstepwidestar[4]{w\in U}{\rA}
  \proofstepwidestar[4]{w\in\FiniteWords A\land P(w)}{%
    \FormulaRefAuto{x\in\{u\in A\mid P(u)\}\eqvdash x\in A\land P(x)}[thm]{4}}
  \proofstepwidestar[4]{w\in\FiniteWords A}{\rAEa{5}}
  \proofstepwidestar[4]{P(w)}{\rAEb{5}}
  \proofstepwidestar[4]{w\in\FiniteOccurrenceSets A}{\FormulaRefAuto{FiniteWordInUniverse}[thm]{6}}
  \proofstepwidestar[]{U\subseteq\FiniteOccurrenceSets A}{%
    \FormulaRefAuto{SubsetIntroductionRule}[thm]{[4]\vdots 8}}
  \proofstepwidestar[]{\varnothing\in\FiniteWords A}{\FormulaRefAuto{GeneratedWordEmpty}[thm]}
  \proofstepwidestar[1]{\varnothing\in U}{%
    \FormulaRefAuto{x\in\{u\in A\mid P(u)\}\eqvdash x\in A\land P(x)}[thm]{\rAI{10,1}}}
  \proofstepwidestar[12]{a\in A}{\rA}
  \proofstepwidestar[2,4]{\forall a\in A\,
    \bigl(P(w)\rightarrow P(\OccurrenceAdjunction w a)\bigr)}{\rRE{6,\rUE{2}}}
  \proofstepwidestar[2,4,12]{P(w)\rightarrow P(\OccurrenceAdjunction w a)}{\rRE{12,\rUE{13}}}
  \proofstepwidestar[2,4,12]{P(\OccurrenceAdjunction w a)}{\rRE{7,14}}
  \proofstepwidestar[4,12]{\OccurrenceAdjunction w a\in\FiniteWords A}{%
    \FormulaRefAuto{GeneratedWordClosed}[thm]{6,12}}
  \proofstepwidestar[2,4,12]{\OccurrenceAdjunction w a\in U}{%
    \FormulaRefAuto{x\in\{u\in A\mid P(u)\}\eqvdash x\in A\land P(x)}[thm]{\rAI{16,15}}}
  \proofstepwidestar[2]{\forall w\in U\,\forall a\in A\,(\OccurrenceAdjunction w a\in U)}{\rUI{\rRI{4,\rUI{\rRI{12,17}}}}}
  \proofstepwidestar[1,2]{\FiniteWords A\subseteq U}{%
    \FormulaRefAuto{GeneratedWordMinimality}[thm]{9,11,18}}
  \proofstepwidestar[1,2,3]{v\in U}{%
    \FormulaRefAuto{A\subseteq B\dsep x\in A\vdash x\in B}[thm]{19,3}}
  \proofstepwidestar[1,2,3]{v\in\FiniteWords A\land P(v)}{%
    \FormulaRefAuto{x\in\{u\in A\mid P(u)\}\eqvdash x\in A\land P(x)}[thm]{20}}
  \proofstepwidestar[1,2,3]{P(v)}{\rAEb{21}}
\end{tabproofwide}


\Needspace{22\baselineskip}
\noindent\begin{minipage}{\linewidth}
\FormulaThmDeltaK[Induktionsregel für erzeugte Wörter]{%
  \begin{aligned}[t]
    &P(\varnothing)\dsep{}\\[-2pt]
    &[u\in\FiniteWords A,a\in A,P(u)]\vdots
      P(\OccurrenceAdjunction u a)\dsep{}\\[-2pt]
    &v\in\FiniteWords A\vdash P(v)
  \end{aligned}
}{GeneratedWordInductionRule}{%
  \DeltaRow{Mengen}{A\dsep u\dsep a\dsep v}
  \DeltaRow{Einstelliges Prädikat}{P}
}
\end{minipage}\par
Die Regel fasst Anfang, lokalen Anfügungsschritt und Induktionsschluss
zusammen, wie die Induktionsregel aus Band~10. Die Schrittvariablen
\(u,a\) sind beliebig; sie dürfen in den übrigen offenen Annahmen nicht
frei vorkommen. Zusätzliche Parameter des Prädikats bleiben fest.
Eine lokale Annahme darf im Schritt unbenutzt bleiben.
\Needspace{9\baselineskip}
\begin{tabproofwide}
  \proofstepwidestar[1]{P(\varnothing)}{\rA}
  \proofstepwidestar[2]{v\in\FiniteWords A}{\rA}
  \proofstepwidestar[3]{u\in\FiniteWords A}{\rA}
  \proofstepwidestar[4]{a\in A}{\rA}
  \proofstepwidestar[5]{P(u)}{\rA}
  \proofstepwidestar[3,4,5]{P(\OccurrenceAdjunction u a)}{%
    \begin{aligned}[t]
      &\text{Schrittprämisse}\\[-2pt]
      &(3,4,5)
    \end{aligned}}
  \proofstepwidestar[]{\forall u\in\FiniteWords A\,\forall a\in A\,
    (P(u)\rightarrow P(\OccurrenceAdjunction u a))}{%
    \rUI{\rRI{3,\rUI{\rRI{4,\rRI{5,6}}}}}}
  \proofstepwidestar[1,2]{P(v)}{%
    \FormulaRefAuto{GeneratedWordInduction}[thm]{1,7,2}}
\end{tabproofwide}

In Anwendungen genügen die Verweise auf Induktionsanfang (IA),
Induktionsschritt (IS) und Wortzugehörigkeit. Bei einem Schritt innerhalb
derselben Tabelle bezeichnet \([i,j,k]\vdots l\) die Teilableitung mit
den lokalen Annahmen \(i,j,k\) und der Endzeile \(l\). Diese Annahmen
werden beim Induktionsschluss entlassen. Für einen natürlichen Index
\(n\) steht bereits \FormulaRefAuto{PeanoInductionRuleAddOne}[thm] zur
Verfügung; dabei wird \(P(n)\) über Zahlen statt über Wörter bewiesen.
